\section{Základy kryptografie (BPC-ZKR)}
    \subsection{Prvočísla – generování, testování, základní testy, typy prvočísel}
        \subsubsection{Prvočísla}
            \begin{itemize}
                \item Prvočíslo je definováno jako celé číslo $\geq 2$, které má pouze triviální dělitele
                \item Prvočísel existuje nekonečně mnoho a jejich počet (hustotu) menší než dané číslo $n$ lze aproximovat vztahem $\pi(n) \approx \frac{n}{\ln(n)}$. 
                \item V moderní kryptografii se běžně pracuje s prvočísly o velikosti 1024 až 2048 bitů (308 až 617 cifer)
            \end{itemize}
        \subsubsection{Typy prvočísel}
            Pro kryptografické účely se často nevyužívají obecná prvočísla, ale prvočísla ve specifickém tvaru, která zajišťují požadované vlastnosti:
            \begin{itemize}
                \item \textbf{Mersennova prvočísla ($M_p$):} Mají tvar $M_p = 2^p - 1$, přičemž exponent $p$ musí být rovněž prvočíslo.
                \item \textbf{Sophie-Germain prvočísla ($SG_p$):} Jedná se o prvočíslo $SG_p$, pro které platí, že i číslo $p = 2SG_p + 1$ je prvočíslem.
                \item \textbf{Bezpečná prvočísla (Safe primes - $B_p$):} Mají tvar $B_p = 2p + 1$, kde $p$ je prvočíslo.
                \end{itemize}
        \subsubsection{Generování prvočísel}
            \begin{itemize}
                \item \textbf{Eratosthenovo síto:} Historický a jednoduchý způsob. Pro velká čísla je nepraktický. Algoritmus sestaví seznam čísel od 2 do n. První nevyřazené číslo označí za prvočíslo a všechny jeho násobky ze seznamu odstraní. Tento postup se opakuje, dokud se nedosáhne hodnoty $\sqrt{n}$; zbylá čísla v seznamu představují prvočísla.
                \item \textbf{Praktický postup v kryptografii:} Generování velkých prvočísel probíhá iterativně: nejprve se vygeneruje náhodné číslo o požadované velikosti (bitové délce) a následně se otestuje, zda se jedná o prvočíslo. Pokud testem neprojde, celý proces generování a testování se opakuje s novým kandidátem.
            \end{itemize}
        \subsubsection{Testování prvočísel}
            \subsubsection*{Testy pravděpodobné prvočíselnosti}
                Z důvodu výkonu a výpočetní náročnosti se u obrovských čísel používají pravděpodobnostní testy. Tyto testy neodkážou prvočíselnost absolutně prokázat, pravděpodobnost omylu je ale zanedbatelná.
                \begin{itemize}
                    \item \textbf{Fermatův test:} Tento test je založen na Malé Fermatově větě, podle níž pro prvočíslo $n$ a s ním nesoudělné kladné číslo $a$ platí kongruence $a^{n-1} \equiv 1 \pmod{n}$. Postup spočívá v náhodné volbě základu $a$ v intervalu $2 \le a \le n - 2$. Pokud kongruence neplatí, testované číslo je s jistotou složené. Hlavní slabinou testu je selhání na tzv. Carmichaelových číslech -- složených číslech, pro která kongruence platí pro libovolné $a$ nesoudělné s $n$. Z tohoto důvodu se v reálné praxi Fermatův test téměř nepoužívá.
                    
                    \item \textbf{Miller-Rabinův test:} Test, který spolehlivě vylučuje Carmichaelova čísla. Složené číslo jím dokáže projít v jedné iteraci s maximální pravděpodobností pouze 25 \%. Kandidát $n$ se vyjádří ve tvaru $n - 1 = 2^s r$, kde $r$ je liché číslo. Poté se náhodně zvolí základ $a$ ($2 \le a \le n - 2$) a testuje se, zda platí $a^r \equiv 1 \pmod{n}$, nebo zda $a^{2^j r} \equiv -1 \pmod{n}$ pro některé $j$ ($0 \le j \le s - 1$). Pokud neplatí ani jedna z podmínek, základ $a$ je odhalen jako \textit{silný svědek složenosti} a číslo $n$ je složené. Pokud je alespoň jedna podmínka splněna, iterace se zopakuje s novým základem $a$.
                \end{itemize}
            \subsubsection*{Testy skutečné prvočíselnosti}
                \begin{itemize}
                    \item \textbf{Lucas-Lehmerův test:} Tento test je výhradně na testování Mersennových prvočísel ve tvaru $n = 2^s - 1$. Využívá se při něm rekurentní postup s počáteční hodnotou $u_0 = 4$ a následným vzorcem $u_{k+1} = (u_k^2 - 2) \pmod{n}$. Test postupně spočte prvky posloupnosti $u_1, u_2, \dots, u_{s-2}$. Pokud je finálním výsledkem $u_{s-2} = 0$, je nepochybně matematicky dokázáno, že se jedná o Mersennovo prvočíslo.
                \end{itemize}
                
        \subsection{Klasifikace výpočetních problémů - složitost algoritmů, třídy složitosti, popis problémů používaných v kryptografii (RSA, diskrétní logaritmus, DH)}
            \subsubsection{Složitost algoritmů}
                Teorie složitosti klasifikuje výpočetní problémy podle jejich náročnosti na tyto prostředky. Základní metriky složitosti jsou \textbf{Časová složitost} (vyjadřuje vztah mezi velikostí vstupu a počtem operací potřebných k vykonání algoritmu) a \textbf{Paměťová složitost} (vyjadřuje vztah mezi velikostí vstupu a paměťovými nároky potřebnými pro běh algoritmu).
                \subsubsection*{Asymptotická složitost}
                Přesné určení složitosti je často nemožné, využívá se asymptotické značení pro různé stavy běhu algoritmu:
                \begin{itemize}
                    \item \textbf{Notace $\mathcal{O}(f(n))$ (Omikron):} Ohraničuje funkci shora (nejhorší možný případ)
                    \item \textbf{Notace $\Omega(f(n))$ (Omega):} Ohraničuje funkci zespodu. (nejlepší možný případ)
                    \item \textbf{Notace $\Theta(f(n))$ (Theta):} Označuje složitost, která má stejný řád jako funkce $f(n)$. ('střední' případ)
                \end{itemize}
                Základní asymptotické třídy (seřazené dle rychlosti růstu): $\mathcal{O}(1)$ (konstantní), $\mathcal{O}(\log \log n)$ (dvojitě logaritmická), $\mathcal{O}(n)$ (lineární), $\mathcal{O}(n \log n)$ (lineárně-logaritmická), $\mathcal{O}(n^2)$ (kvadratická), $\mathcal{O}(n^3)$ (kubická), $\mathcal{O}(n^p)$ (polynomiální), $\mathcal{O}(p^n)$ (exponenciální) a $\mathcal{O}(n!)$ (faktoriálová).

                \subsubsection{Třídy složitosti}
                    Pro klasifikaci výpočetní náročnosti problémů se definují základní třídy složitosti:
                    \begin{itemize}
                        \item \textbf{Třída P (Polynomiální):} Obsahuje všechny problémy, které dokáže deterministický algoritmus \textit{vyřešit} v polynomiálním čase.
                        \item \textbf{Třída NP (Nedeterministicky polynomiální):} Obsahuje problémy, pro které sice nemusí existovat algoritmus řešící je v polynomiálním čase, ale pokud je předloženo jejich řešení, lze v polynomiálním čase \textit{ověřit} jeho správnost.
                        \item \textbf{Třída NP-úplná:} Představuje množinu 'nejtěžších' problémů ve třídě NP. Významná vlastnost NP-úplných problémů je, že jakýkoliv jiný problém ze třídy NP lze na ně v polynomiálním čase redukovat.
                    \end{itemize}
                    Vztah mezi třídami není matematicky zcela dokázán. V informatice a kryptografii se všeobecně předpokládá platnost tvrzení $P \neq NP$. Z toho plyne, že těžké problémy typu faktorizace se nedají v rozumném čase vyřešit pouhou hrubou silou nebo současnými algoritmy -- využívá se k budování asymetrických kryptosystémů.

                \subsubsection{Popis problémů používaných v kryptografii}
                    \begin{itemize}
                        \item \textbf{Problém faktorizace (RSA):} Kryptosystém RSA je založen na modulárních operacích, kde modulem je složené číslo $n = r \cdot s$, přičemž $r$ a $s$ jsou velká prvočísla. Bezpečnost RSA spočívá výhradně v obtížnosti problému faktorizace velkých čísel. Útočník potřebuje soukromý klíč, jehož výpočet závisí na Eulerově funkci $\varphi(n) = (r - 1)(s - 1)$. Pro jeho získání je nutné veřejně známé velké číslo $n$ (v praxi 1024--2048 bitů) rozložit zpět na prvočinitele $r$ a $s$. Pro klasické výpočetní systémy je nalezení takového rozkladu výpočetně neřešitelné.
                        
                        \item \textbf{Problém diskrétního logaritmu (Diffie-Hellman a DSA):} Algoritmy využívající modulární umocňování na multiplikativních grupách vycházejí z problému diskrétního logaritmu. V protokolu DH strany bezpečně ustaví sdílený symetrický klíč nad velkým prvočíselným modulem $p$ pomocí generátoru $g$ tak, že strana vypočte a rozešle $A = g^a \pmod{p}$. Modulární umocňování se provede rychle, avšak inverzní extrakce exponentu $a$ ze znalosti $A, g$ a $p$ (výpočet diskrétního logaritmu) je výpočetně složitá.
                        
                        \item \textbf{Diskrétní logaritmus nad eliptickými křivkami (EC-DH / EC-DSA):} Jako alternativní algebraické struktury se používají eliptické křivky (ECC). Místo modulárního umocňování se používá násobení bodu křivky skalárem (celým číslem), což se provádí jako $a$-násobné přičtení bodu k sobě samému. Inverzní operace (zjištění původního skaláru $a$) je výpočetně ještě náročnějším problémem než tradiční diskrétní logaritmus nad okruhem celých čísel.
                    \end{itemize}
            \subsection{Symetrické šifrování – princip, využití, algoritmy DES a AES. Hashovací funkce – požadavky, zástupci }
                \subsubsection{Symetrické šifrování}
                    Symetrické kryptosystémy jsou založeny na předpokladu, že obě komunikující strany sdílejí totožný tajný klíč, případně klíče, které lze jeden z druhého velmi snadno odvodit. Podle způsobu zpracování vstupních dat se tyto systémy dělí na \textbf{proudové šifry}, které zpracovávají data bit po bitu nebo bajt po bajtu, a na \textbf{blokové šifry}, které zpracovávají data po větších, fixně daných blocích. 

                    Symetrické šifrování se primárně nasazuje při šifrování velkých objemů dat. Mezi jeho zásadní výhody patří výpočetní nenáročnost, rychlost a efektivní implementovatelnost do hardwaru. Kryptografickou slabinou symetrických systémů je problematika bezpečné distribuce tajných klíčů. Tyto systémy zároveň samy o sobě neřeší otázku nepopiratelnosti.
                \subsubsection{Algoritmy DES a AES}
                    Při konstrukci blokových šifer se uplatňují dva základní návrhové vzory: Feistelovo schéma (vstup se dělí na poloviny, které se s klíčem křížově promíchávají) a SP-síť (Substituční-permutační síť).
                    \begin{itemize}
                        \item \textbf{DES (Data Encryption Standard)}
                            \begin{itemize}
                                \item Šifra je postavená na Feistelově schématu, jež zpracovává data po 64bitových blocích. Šifra využívá klíč, který byl uměle oslaben na 56 bitů.
                            \end{itemize}
                            \begin{figure}[!h]
                                \begin{center}
                                \includegraphics[scale=0.4]{Feistel.png}
                                \caption{První část Feistelova schématu}
                                \end{center}
                            \end{figure}\textbf{}
                        \item \textbf{AES (Advanced Encryption Standard)}
                            \begin{itemize}
                                \item AES je bloková šifra navržená na bázi SP-sítě. Standardizovaný vstupní blok má délku 128 bitů, pro klíč jsou povoleny délky 128, 192, nebo 256 bitů.
                                \item Šifrování probíhá v několika iteračních rundách, jejichž počet závisí na délce klíče. Každá standardní runda aplikuje na blok čtyři hardwarově nenáročné operace: \textit{ByteSub} (nelineární substituce jednotlivých bajtů přes S-box), \textit{ShiftRow} (cyklický posun bajtů v řádcích), \textit{MixColumn} (lineární kombinování bajtů ve sloupcích zajišťující difúzi) a \textit{AddRoundKey} (kombinace bloku s rundovním klíčem pomocí operace XOR).
                            \end{itemize}
                    \end{itemize}
                \subsubsection{Hashovací funkce – požadavky, zástupci}
                    Hashovací funkce je matematický jednosměrný algoritmus, který pro libovolně velký vstupní datový soubor vypočítá otisk (tzv. hash či digest) reprezentovaný řetězcem o pevně dané délce.

                    Mezi historické hašovací funkce patří například SHA-1 nebo MD5, jehož hash je zranitelný vůči kolizím. V současnosti se používají algoritmy SHA-2 (typicky varianty s výstupem 224, 256, 384 a 512 bitů) nebo jeho nástupce SHA-3.

                    
                    \subsubsection*{Požadavky na bezpečnost:}
                        \begin{enumerate}
                            \item \textbf{Jednosměrnost:} Z jakéhokoliv vstupu $m$ je výpočetně velmi jednoduché získat hash $H(m)$, avšak opačný proces – zrekonstruovat původní data ze znalosti samotného hashe – musí být výpočetně neřešitelný.
                            \item \textbf{Bezkoliznost 1. řádu:} Je výpočetně prakticky nemožné objevit libovolné dvě zcela odlišné zprávy $m$ a $m'$, jejichž výstupní hashe by byly totožné ($H(m) = H(m')$).
                            \item \textbf{Bezkoliznost 2. řádu:} Pro jednu konkrétní fixní zprávu $m$ je výpočetně nemožné nalézt modifikovanou (podvrženou) zprávu $m'$, která by po zahašování dávala naprosto stejný výsledek ($H(m) = H(m')$).
                            \item \textbf{Rychlost:} Hashovací funkce musí být rychlá
                        \end{enumerate}
            \subsection{Asymetrické šifrování – princip, využití, algoritmus RSA, protokol Diffie-Hellman}
                Asymetrická kryptografie využívá dvojici matematicky svázaných klíčů ($K \neq K'$). Tuto dvojici tvoří veřejný klíč ($Pk$), který je volně k dispozici, a soukromý klíč ($Sk$), který je držen v tajnosti. Proces šifrování probíhá pomocí veřejného klíče, dešifrování zase pomocí soukromého.
                
                Asymetrické šifrování řeší problém distribuce klíčů, se kterým se potýkají symetrické systémy. Nevýhodou asymetrických systémů je jejich vysoká výpočetní náročnost ve srovnání se symetrickými algoritmy.
                \subsubsection{Algoritmus RSA (Rivest, Shamir, Adleman)}
                    Bezpečnost RSA je postavena na problému faktorizace. Systém pracuje s modulem $n$, který je součinem dvou velkých prvočísel ($n = r \cdot s$). Pro zajištění bezpečnosti se vyžaduje, aby délka modulu $n$ činila minimálně 2048 bitů.
                    \subsubsection*{Generování parametrů, proces šifrování a dešifrování:}
                        \begin{enumerate}
                            \item Zvolí se dvě velká bezpečná prvočísla $r$ a $s$.
                            \item Vypočte se modulus $n = r \cdot s$.
                            \item Vypočte se hodnota Eulerovy funkce $\varphi(n) = (r - 1)(s - 1)$.
                            \item Zvolí se soukromý klíč $Sk$ tak, aby byl s $\varphi(n)$ nesoudělný, tj. $GCD(Sk, \varphi(n)) = 1$.
                            \item Vypočte se veřejný klíč $Pk \equiv Sk^{-1} \pmod{\varphi(n)}$.
                            \item Kryptogram $c = m^{Pk} \pmod{n}$, Zpráva $m = c^{Sk} \pmod{n}$.
                        \end{enumerate}
                \subsubsection{Protokol Diffie-Hellman (DH)}
                    Diffie-Hellmanův protokol slouží k bezpečnému ustanovení sdíleného tajného klíče přes nedůvěryhodné komunikační médium, je však náchylný k útoku Man-in-the-Middle. Bezpečnost protokolu DH se opírá o výpočetní složitost problému diskrétního logaritmu v multiplikativní grupě ${Z}_p^*$. Veřejnými předem sdílenými parametry jsou velké bezpečné prvočíslo $p$ (modulus o délce alespoň 3072 bitů) a číslo $g$, což je generátor podgrupy řádu $q$ (délka alespoň 256 bitů).

                    \subsubsection*{Průběh protokolu:}
                        \begin{enumerate}
                            \item \textbf{Účastník 1:} Zvolí si zcela náhodné číslo $a \in {Z}_q$. Spočte hodnotu $A = g^a \pmod{p}$ a veřejně ji odešle.
                            \item \textbf{Účastník 2:} Nezávisle zvolí náhodné číslo $b \in {Z}_q$. Spočte hodnotu $B = g^b \pmod{p}$ a veřejně ji odešle.
                            \item \textbf{Výpočet klíče (Účastník 1):} Přijatou hodnotu $B$ umocní svým soukromým klíčem $a$, čímž získá výsledný symetrický klíč $K = B^a \pmod{p}$.
                            \item \textbf{Výpočet klíče (Účastník 2):} Přijatou hodnotu $A$ umocní svým soukromým klíčem $b$, čímž získá tentýž symetrický klíč $K = A^b \pmod{p}$. 
                        \end{enumerate}
                    
            \subsection{Autentizace – princip protokolů výzva-odpověď, certifikáty, digitální podpis.}
                Autentizace je proces ověření pravosti. V kryptografii a bezpečnosti se dělí na dvě hlavní oblasti:

                \begin{itemize}
                    \item \textbf{Autentizace dat:} Zajišťuje ověření, že data skutečně pocházejí od očekávaného odesílatele a nacházejí se v původní formě. K tomuto účelu se v asymetrické kryptografii využívá elektronický (digitální) podpis.
                    \item \textbf{Autentizace entit (osob/zařízení):} Zajišťuje spolehlivé ověření, že komunikujeme se zamýšleným subjektem.
                \end{itemize}

                \subsubsection{Princip protokolů Challenge-Response}
                     K bezpečné autentizaci osob (bez nutnosti přenášet heslo v otevřené podobě) slouží protokoly postavené na mechanismu \textbf{challenge-response}.
                    
                    \begin{itemize}
                        \item \textbf{Princip fungování:} Server zašle klientovi náhodnou \textit{výzvu} (challenge). Klient ze znalosti svého tajného hesla a přijaté výzvy vypočte kryptografickou \textit{odpověď} (např. zašifrováním výzvy nebo vytvořením otisku pomocí hashovací funkce) a tu odešle zpět serveru. Server provede na své straně totožný výpočet a ověří, zda se jím vypočtená hodnota shoduje s přijatou odpovědí ($\textit{odpoved} = F(\textit{vyzva}, \textit{heslo})$).
                        \item \textbf{Praktické nasazení:} Tento princip nachází uplatnění například v rozšířené autentizaci CHAP v rámci protokolu RADIUS nebo doménové autentizace NTLM v systémech Windows, kde server volí náhodnou výzvu a klient na ni odpovídá derivátem svého hashe hesla.
                    \end{itemize}
                \subsubsection{Digitální podpis}
                    Digitální podpis slouží k jasnému prokázání autorství a integrity odeslané zprávy pomocí asymetrické kryptografie. Podepisování a ověřování funguje na základě páru matematicky svázaných klíčů -- soukromého a veřejného.
                    
                    \begin{itemize}
                        \item \textbf{Postup podepisování a ověření:} Odesílatel zkrátí původní zprávu na její otisk pomocí hashovací funkce. Tento otisk následně zašifruje svým drženým \textbf{soukromým klíčem} ($Sk$), čímž vznikne digitální podpis ($Sig_m$). Příjemce použije odpovídající \textbf{veřejný klíč} ($Pk$) odesílatele. Pomocí něj dešifruje přijatý podpis a výsledek porovná se spočteným hashem přijaté zprávy.
                        \item \textbf{Využívané algoritmy:} Standardním zástupcem je kryptosystém \textbf{RSA}, u kterého matematické operace digitálního podpisu pracují s modulem tvořeným součinem dvou velkých prvočísel. Dalším významným zástupcem je \textbf{DSA} (Digital Signature Algorithm), jež operuje nad modulární multiplikativní grupou modulo prvočíslo a pro každý jednotlivý podpis vyžaduje vygenerování nové, unikátní náhodné hodnoty $k$.
                    \end{itemize}
                \subsection*{4. Certifikáty}
                    Certifikáty se využívají k ověření autentičnosti osob a koncových zařízení. Jejich hlavním smyslem je kryptograficky svázat konkrétní identitu uživatele nebo serveru s jeho veřejným klíčem, což brání útokům typu Man-in-the-Middle.
                    
                    \begin{itemize}
                        \item \textbf{Princip fungování v SSL/TLS:} Při sestavování zabezpečeného spojení odešle server klientovi svůj certifikát ($Cert_S$), který obsahuje veřejný klíč serveru. Klient tento certifikát a jeho validitu ověří. Následně může bezpečně vytvořit sdílené tajemství, zašifrovat jej veřejným klíčem vytaženým z certifikátu a odeslat zpět serveru, z čehož obě strany finálně odvodí symetrické šifrovací klíče. V případě potřeby oboustranného ověření je nasazena tzv. vzájemná autentizace, kde svůj vlastní certifikát ($Cert_C$) navíc předkládá k ověření i klient.
                    \end{itemize}
\newpage